\section{Local \texorpdfstring{$\aleph_1$}{א{1}}-presentability of \texorpdfstring{\cstars{}s}{σC*-algebras}}\label{sec:accessible}

In this section, we ensure local $\aleph_1$-presentability of $\SigmaC$ and accessibility of the forgetful functor $U$.
The reader is advised to have some background on presentable and accessible categories (see, e.g., \cite[Section 1 and 2]{adamek1994locallypresentable} for some definitions that will be used below).
For the sake of readability, we recall that a poset $P$ is \newterm{$\aleph_1$-directed} if every countable subset has an upper bound.

We start with a technical result.

\begin{lemma}\label{lem:technical_access}
	Let us consider an $\aleph_1$-directed diagram $D\colon P \to \Calg$, by which we mean a functor whose domain $P$ is an $\aleph_1$-directed poset. Call $\A\coloneqq \colim D$, and $\phi_{\alpha',\alpha}\colon D(\alpha')\to D(\alpha)$ the transition morphisms, and $\phi_{\alpha}\colon D(\alpha)\to \A$ the maps of the initial cocone. Then,
	\begin{enumerate}
		\item\label{it:technical_access_norm}\emph{Isometric representation}. For every $a\in \A$, there exist $\alpha \in P$ and $a_{\alpha}\in D(\alpha)$ such that $a=\phi_{\alpha}(a_{\alpha})$.
		      Moreover, $a_{\alpha}$ can be chosen so that $\norm{a_{\alpha}}_{D(\alpha)} = \norm{a}_{\A}$.
		\item\label{it:technical_access_positivity}\emph{Positive representation}. For every positive $a\in \A$, there exists $\alpha\in P$ and a positive $a_{\alpha}\in D(\alpha)$ such that $a=\phi_{\alpha}(a_{\alpha})$.
	\end{enumerate}
\end{lemma}
\begin{proof}\item
	\begin{enumerate}
		\item[\ref{it:technical_access_norm}]
		      Consider the quotient $B=\bigcup_{\alpha \in P} D(\alpha)/_\sim$, i.e.\@\xspace the colimit of the diagram $D$ in the category of $*$-algebras.
		      As discussed in \cite[Corollaries 7.22 and 7.23]{bunke2021lectures}, $\A$ corresponds to the norm-completion of $B$. %(a fact we already used in \cref{prop:tensor_colims}).     

		      We claim that for $a\in B$, we can find $\alpha$ and $a_{\alpha}\in D(\alpha)$ such that $\phi_{\alpha}(a_{\alpha})=a$ and $\norm{a_{\alpha}}_{D(\alpha)}=\norm{a}_B$.
		      Recall that the norm is defined as
		      \[\norm{a}_B \coloneqq \inf \norm{a_{\alpha}}_{D(\alpha)},\]
		      where the infimum is calculated over all $\alpha \in P$ and $a_{\alpha}\in D(\alpha)$ such that $\phi_{\alpha}(a_{\alpha})=a$.
		      In particular, for each $n$, we can find $\alpha_n$ and $a_n\in D(\alpha_n)$ such that $\phi_{\alpha_n}(a_n)=a$ and $\norm{a_n}_{D(\alpha_n)} \le \norm{a}_B+ \frac{1}{n}$.
		      By definition of $B$, $a_n \sim a_{m}$ for all $n,m$.
		      Since $P$ is directed, this means that there exists a certain $\alpha>\alpha_n,\alpha_m$ such that $\phi_{\alpha_n,\alpha} (a_n) = \phi_{\alpha_m,\alpha}(a_m)$.
		      By letting $m$ vary, since every countable subset of $P$ has an upper bound, such an $\alpha$ can be chosen to identify all the elements $a_n$, i.e.\  $\phi_{\alpha_n,\alpha} (a_n) = \phi_{\alpha_m,\alpha}(a_m)$ for all $m,n$ simultaneously.
		      We now define $a_{\alpha}\coloneqq \phi_{\alpha_n,\alpha} (a_n)\in D(\alpha)$, which does not depend on $n$ by the argument above.
		      Since $\phi_{\alpha_n,\alpha}$ is a $*$-homomorphism, it is contractive, and therefore
		      \[ \norm{a_{\alpha}}_{D(\alpha)}\le \norm{a_n}_{D(\alpha_n)}\le \norm{a}_B+ \frac{1}{n} \]
		      from which $\norm{a_{\alpha}}_{D(\alpha)}= \norm{a}_B$ by arbitrariness of $n$.
		      Finally, the claim is proved.


		      We now seek to show that $B=\A$.
		      Since $\A$ corresponds to the norm-completion of $B$ (as mentioned at the beginning of the proof), it suffices to take any Cauchy sequence $(a_n)$ of $B$ and show that it admits a limit in $B$.

		      Passing to a subsequence if necessary, we may assume that $\norm{a_n-a_{n-1}}\le \frac{1}{2^{n}}$ for all $n\ge 2$.
		      For each difference $d_n\coloneqq a_n-a_{n-1}$, with $n\ge 1$ (we set $a_0\coloneqq 0$), one can find $\alpha_n$ and $d'_n$ such that $\phi_{\alpha_n}(d'_n)=d_n$ and $\norm{d'_n}=\norm{d_n}$, hence $\norm{d'_n}\le \frac{1}{2^{n}}$.
		      By taking the supremum $\alpha$ of all $\alpha_n$, we can assume that $(d'_n)\subseteq D(\alpha)$.
		      In other words, we identify $d'_n$ with its image in $D(\alpha)$.
		      We note that the norms are not increased by this assumption, since all transition morphisms $\phi_{\alpha_n,\alpha}$ are contractive.
		      Then, the sequence $s_n \coloneqq \sum_{i=1}^n d'_i$ is Cauchy in $D(\alpha)$ by direct check: for $m>n\ge 1$,
		      \[
			      \norm{s_m - s_n} =\norm{\sum_{i=n+1}^m d_i} \le \sum_{i=n+1}^m \frac{1}{2^i} = \frac{1}{2^n} \left( \sum_{i=1}^{m-n} \frac{1}{2^i}\right) \le \frac{1}{2^n} \xrightarrow{n,m\to \infty} 0.
		      \]
		      Moreover, its limit $a_{\alpha}$ in $D(\alpha)$ yields a limit for $(a_n)$ because $\phi_{\alpha}$ is bounded:
		      \[
			      \phi_{\alpha}(a_{\alpha}) =\lim_n \phi_{\alpha}(s_n) = \lim_n \left(\sum_{i=1}^n d_i \right)= \lim_n a_n.
		      \]
		      Finally, $B=\A$, so the statement follows by the claim we proved above.
		\item[\ref{it:technical_access_positivity}] Let $a\in \A$ be a positive element.
		      By Item~\ref{it:technical_access_norm}, there exists some $\alpha$ and $a'_{\alpha}$ satisfying $\phi_{\alpha}(a'_{\alpha})=a$.
		      We now take $a_{\alpha}\coloneqq \max(a'_{\alpha},0)\ge 0$, where this notation is a standard application of the continuous functional calculus (see, e.g., \cite[Proposition 4.6]{takesaki1979operatoralgebras1}).
		      By functoriality of functional calculus~\cite[Corollary 8.4]{takesaki1979operatoralgebras1},
		      \[
			      \phi_{\alpha}(a_{\alpha})= \phi_{\alpha}(\max(a'_{\alpha},0))=\max(\phi_{\alpha}(a'_{\alpha}),0)=\max(a,0) = a
		      \]
		      where the last equality follows from $a\ge 0$.\qedhere
	\end{enumerate}
\end{proof}

\begin{remark}\label{rem:Calg_accessible}
	We recall that $\Calg$ is a locally $\aleph_1$-presentable category (combine \cite[Theorem 2.4]{pelletier1989equationalCstar} and \cite[Theorem 3.28]{adamek1994locallypresentable}). %\url{https://arxiv.org/pdf/2203.12087}. 
    The result of Pelletier and Rosick\'y quoted here is the closest antecedent of the present programme: it exhibits \cstar{}s as the models of an equational theory. 
	We establish a weaker analogue for \cstars{}s, namely that they are the models of an $\aleph_1$-ary essentially algebraic theory (\cref{cor:essentially_algebraic}). 
	Whether they are also the models of an equational theory is left open.
\end{remark}

\begin{theorem}\label{thm:U_accessible}
	The functor $U\colon \SigmaC \to \Calg$ is $\aleph_1$-accessible and $\SigmaC$ is locally $\aleph_1$-presentable.
\end{theorem}
\begin{proof}
	Let us consider an $\aleph_1$-directed diagram $D\colon P \to \SigmaC$, where $P$ is an $\aleph_1$-directed poset.
	We will show that $U(\colim D)\cong \colim UD$. 
	By \cite[Theorem 1.5 and Remark 1.21]{adamek1994locallypresentable} this will be enough to ensure a similar result holds for all colimits over $\aleph_1$-\emph{filtered} categories, not necessarily posets. 

	Concretely, since $U$ simply forgets $\sigma$-normality, it suffices to ensure that $\A \coloneqq \colim UD$ is a \cstars{} and the cocone morphisms are $\sigma$-normal.
	Once this is established, since left adjoints always preserve colimits, it is clear that $UF=T\colon \Calg \to \Calg$ is an $\aleph_1$-accessible monad over a locally $\aleph_1$-presentable category (see \cref{rem:Calg_accessible}), so $\SigmaC$ is also locally $\aleph_1$-presentable by combining \cref{thm:U_monadic} and \cite[Theorem and Remark 2.78]{adamek1994locallypresentable}.
	Finally, $U$ is indeed $\aleph_1$-accessible as claimed.

	By \cref{lem:technical_access}\ref{it:technical_access_norm}, all elements $a\in \A\coloneqq\colim UD$ are given by some $a_{\alpha}\in D(\alpha)$, where $\alpha$ is arbitrary.
	From now on, let $(a_n)$ be a norm-bounded monotone increasing sequence in $\A$. %We aim to show that $\sup_n a_n$ exists in $\A$.
	Below, we use $\phi_{\alpha',\alpha}$ and $\phi_{\alpha}$ to denote the transition morphisms and the cocone morphisms, respectively, as in the statement of \cref{lem:technical_access}.

	We first show that there exists $\gamma \in P$ and a norm-bounded monotone increasing sequence $(b_n)\subseteq D(\gamma)$ satisfying $\phi_{\gamma}(b_n)=a_n$.
	Let us take $\alpha_1\in P$ such that $a_1$ has a representative $d_1\in D(\alpha_1)$.
	Next, for $n>1$, note that $a_n-a_{n-1}\ge 0$ by assumption.
	Then there exists $\alpha_n\in P$ and a positive representative $d_n \in D(\alpha_n)$ of ${a_n - a_{n-1}}$ by \cref{lem:technical_access}\ref{it:technical_access_positivity}.

	By taking the supremum $\alpha_{\bbN}$ of $(\alpha_n)_n$, $a_n$ has a representative $b_n$ in $D(\alpha_{\bbN})$ given by $b_n\coloneqq \sum_{j=1}^n d_n$.
	By definition, $b_n \le b_{n+1}$.
	Moreover, by \cref{lem:technical_access}\ref{it:technical_access_norm}, for each $a_n$ one can find $\beta_n$ and $c_n \in D(\beta_n)$ such that $\phi_{\beta_n}(c_n) =a_n$ and $\norm{c_n}=a_n$.
	Since $c_n \sim b_n$, we can find a certain $\gamma_n > \alpha_{\bbN} , \beta_n $ such that $\phi_{\alpha_{\bbN},\gamma_n}(b_n) = \phi_{\beta_n, \gamma_n}(c_n)$; in particular, $\norm{\phi_{\alpha_{\bbN},\gamma_n}(b_n)}=\norm{a_n}$ because the transition morphisms are contractive.
	By now taking the upper bound $\gamma$ of all $\gamma_n$, and by an abuse of notation calling $b_n$ the elements $\phi_{\alpha_n ,\gamma}(b_n)\in D(\gamma)$, we can finally ensure that $(b_n)\subseteq D(\gamma)$ is a norm-bounded monotone increasing sequence.


	Next, given $x \ge a_n$ and a norm-bounded monotone increasing sequence $(b_n)\subseteq D(\gamma)$ satisfying $\phi_{\gamma}(b_n)=a_n$,
	we claim that there exists $\alpha>\gamma$ and a representative $y\in D(\alpha)$ of $x$ such that $\phi_{\gamma,\alpha}(\sup_n b_n)\le y$.
	This immediately ensures that $\phi_{\gamma}(\sup_n b_n)$ is indeed the supremum of $a_n$ in $\A$, so $\A$ is a \cstars{}, and moreover the arbitrariness of $(b_n)$ implies that the cocone morphisms are $\sigma$-normal.
	Therefore, this claim suffices to conclude the proof of the statement.

	We are now left to prove the claim of the previous paragraph. By \cref{lem:technical_access}\ref{it:technical_access_positivity}, $x-a_n\ge 0$ has a positive representation in some $D(\delta_n)$. By taking the upper bound $\delta$ of $(\delta_n)_n$ and $\gamma$,
	we have a representative $y_n\ge 0$ of $x-a_n$ in $D(\delta)$ for all $n$.
	Since $\delta\ge \gamma$, $y_n + \phi_{\gamma,\delta}(b_n) \sim y_{m}+ \phi_{\gamma,\delta}(b_m)$ because they are both sent to $x$ via the cocone morphism.
	By letting $m$ vary, we can find some $\alpha\ge \delta$ such that $y\coloneqq \phi_{\delta,\alpha}(y_n + \phi_{\gamma,\delta}(b_n)) = \phi_{\delta, \alpha}(y_{m}+ \phi_{\gamma,\delta}(b_m))$ for all $n,m$.
	Then $y\ge \phi_{\gamma,\alpha}(b_n)$ for all $n$ because $\phi_{\delta,\alpha}(y_n)\ge 0$, and thus $y\ge \sup_n \phi_{\gamma,\alpha}(b_n) = \phi_{\gamma,\alpha}(\sup_n b_n)$, because transition morphisms are $\sigma$-normal.
\end{proof}

\begin{corollary}\label{cor:essentially_algebraic}
	$\SigmaC$ is the category of models of an $\aleph_1$-ary essentially algebraic theory.
\end{corollary}
\begin{proof}
		A category is locally $\lambda$-presentable if and only if it is the category of models of a $\lambda$-ary essentially algebraic theory~\cite[Theorem 3.36]{adamek1994locallypresentable}; it therefore suffices to apply \cref{thm:U_accessible}.
\end{proof}

\begin{remark}\label{rem:presentability_consequences}
\cref{thm:U_accessible} carries further consequences. 
For instance, $\SigmaC$ is wellpowered and co-wellpowered. 
Moreover, the family of Pedersen--Baire envelopes $F\A$ of $\aleph_1$-presentable \cstar{}s $\A$ forms a strong generator of $\SigmaC$, and every \cstars{} is an $\aleph_1$-directed colimit of $\aleph_1$-presentable \cstars{}s.
The latter assertion is the precise form of the informal statement, made in the introduction, that each \cstars{} is built out of a small family of $\aleph_1$-compact ones.
\end{remark}


